% ============================================================================
%  CHƯƠNG 3 — Thu thập và xây dựng tập dữ liệu
% ============================================================================

\section{Thu thập và xây dựng tập dữ liệu}
\label{sec:dataset}

Chương này mô tả chi tiết quy trình thu thập dữ liệu từ các nguồn khác nhau, xây dựng hướng dẫn gán nhãn chuẩn hóa và schema lưu trữ JSONL thống nhất. Đồng thời, chương trình bày chi tiết quy trình lọc dữ liệu từ ảnh thô sang tập dữ liệu sạch, các biện pháp kỹ thuật chống rò rỉ thông tin (data leakage) và thống kê tỷ lệ thiếu nhãn.

% ----------------------------------------------------------------------------
\subsection{Nguồn dữ liệu và quy trình lọc mẫu}
\label{subsec:data-source}

Đồ án kết hợp hai nguồn dữ liệu biên lai tiếng Việt: bộ dữ liệu công khai từ cuộc thi MC-OCR 2021 và bộ dữ liệu tự chụp thực tế từ các cửa hàng bán lẻ khác nhau. Quy trình hình thành tập dữ liệu sạch trải qua các bước lọc nghiêm ngặt được tổng hợp trong \cref{tab:data-filtering-stats}.

\subsubsection{MC-OCR 2021}
Tập dữ liệu gốc từ cuộc thi \gls{mcocr} 2021~\cite{mcocr2021} bao gồm \textbf{2.436 ảnh} biên lai. Tuy nhiên, tập dữ liệu này chứa cả phần public test không đi kèm nhãn ground-truth. Nhóm nghiên cứu chỉ lọc và sử dụng phần dữ liệu huấn luyện có nhãn đầy đủ gồm \textbf{1.094 mẫu}. 

Nhãn gốc của MC-OCR được ánh xạ sang schema thống nhất của đồ án như sau:
\begin{itemize}
  \item \code{SELLER} $\to$ \field{store\_name} (Tên cửa hàng)
  \item \code{SELLER\_ADDRESS} $\to$ \field{address} (Địa chỉ)
  \item \code{TIMESTAMP} $\to$ \field{date} (Ngày giao dịch)
  \item \code{TOTAL\_COST} $\to$ \field{total} (Tổng số tiền thanh toán)
\end{itemize}

\subsubsection{Dữ liệu tự thu thập}
Để tăng tính đa dạng và kiểm định tính thực tiễn, nhóm nghiên cứu tự chụp thêm \textbf{500 ảnh} biên lai thực tế từ siêu thị, nhà hàng, cửa hàng tiện lợi tại Hà Nội. Các hình ảnh được chụp bằng thiết bị di động trong điều kiện ánh sáng, góc chụp biến động mạnh và ảnh có thể bị nghiêng, mờ nhẹ. Việc gán nhãn được thực hiện thủ công bằng công cụ Label Studio\footnote{\url{https://labelstud.io}} để đánh dấu bounding box và nhập giá trị văn bản cho 4 trường thông tin.

\subsubsection{Quy trình lọc trùng lặp ảnh (MD5 Deduplication)}
Nhằm loại bỏ các ảnh trùng lặp hoặc các ảnh chụp cùng một hóa đơn nhiều lần nhưng ở các góc máy gần tương tự nhau, nhóm nghiên cứu chạy thuật toán băm MD5 kiểm tra nội dung ảnh trước khi phân chia dữ liệu:
1. Tính giá trị băm MD5 cho từng tệp ảnh.
2. Nếu phát hiện tệp ảnh có cùng mã MD5 với tệp đã xử lý trước đó, tệp ảnh này sẽ bị loại khỏi tập dữ liệu.
Quy trình này đã phát hiện và loại bỏ \textbf{98 mẫu ảnh trùng lặp} (gồm 63 mẫu từ MC-OCR và 35 mẫu tự chụp).

\begin{table}[htbp]
  \centering
  \caption{Thống kê số lượng mẫu bị loại theo từng nguồn dữ liệu.}
  \label{tab:data-filtering-stats}
  \begin{tabular}{lcccc}
    \toprule
    \textbf{Nguồn dữ liệu} & \textbf{Ảnh thô ban đầu} & \textbf{Loại do thiếu nhãn} & \textbf{Loại do trùng MD5} & \textbf{Mẫu sạch cuối cùng} \\
    \midrule
    MC-OCR 2021 & 2.436 & 1.279 & 63 & 1.094 \\
    Tự thu thập  & 500   & 0     & 35 & 465   \\
    \midrule
    \textbf{Tổng} & \textbf{2.936} & \textbf{1.279} & \textbf{98} & \textbf{1.559} \\
    \bottomrule
  \end{tabular}
\end{table}

% ----------------------------------------------------------------------------
\subsection{Hướng dẫn gán nhãn (Annotation Guideline) và QA/QC}
\label{subsec:annotation-guideline}

Việc gán nhãn dữ liệu tuân thủ nghiêm ngặt bộ quy tắc chuẩn hóa nhằm bảo đảm tính nhất quán cao nhất giữa các nguồn dữ liệu khác nhau:
\begin{itemize}
  \item \field{store\_name}: Chỉ lấy tên thương hiệu/cửa hàng chính thức, loại bỏ các loại hình kinh doanh như "Công ty TNHH", "Cửa hàng bán lẻ" hoặc MST đi kèm.
  \item \field{date}: Chỉ lấy thông tin ngày tháng năm và chuẩn hóa về định dạng YYYY-MM-DD. Loại bỏ các tiền tố như "Ngày:", "Giờ:", hoặc thông tin giờ phút giây.
  \item \field{total}: Chỉ lấy số tiền thanh toán cuối cùng (final payment), loại bỏ đơn vị tiền tệ, dấu phân cách hàng nghìn và các số 0 ở đầu. Không lấy giá trị tạm tính (subtotal) hoặc tiền thừa (change).
  \item \field{address}: Lấy toàn bộ địa chỉ ghi trên biên lai, loại bỏ số điện thoại, website, email đi kèm.
\end{itemize}

\textbf{Quy trình QA/QC (Kiểm soát chất lượng):} 
Sau khi hoàn tất gán nhãn tự chụp, nhóm tiến hành kiểm soát chất lượng ngẫu nhiên trên \textbf{50 mẫu}. Kết quả QA/QC ghi nhận tỷ lệ sai lệch ban đầu là 8.0\% (4 mẫu), chủ yếu liên quan đến việc nhầm lẫn giữa trường \field{total} và subtotal, hoặc gán nhãn địa chỉ bị thiếu dòng. Toàn bộ sai lệch phát hiện đều được nhóm nghiên cứu tiến hành sửa đổi thủ công để đạt độ chính xác nhãn sạch 100\% trước khi chuyển qua bước huấn luyện.

% ============================================================================
\subsection{Schema dữ liệu thống nhất (Unified JSONL)}
\label{subsec:unified-schema}

Toàn bộ dữ liệu sạch từ cả hai nguồn được chuyển đổi về một schema JSONL thống nhất với \textbf{13 trường chính} được mô tả chi tiết trong \cref{tab:jsonl-schema}.

\begin{table}[htbp]
  \centering
  \caption{Mô tả 13 trường dữ liệu trong schema JSONL thống nhất.}
  \label{tab:jsonl-schema}
  \begin{tabularx}{\textwidth}{lL{9.5cm}}
    \toprule
    \textbf{Trường} & \textbf{Mô tả} \\
    \midrule
    \code{id}                & Mã định danh duy nhất của mỗi mẫu ảnh \\
    \code{group\_id}         & Mã nhóm ảnh (dùng để nhận diện các ảnh chụp cùng biên lai) \\
    \code{store\_group}      & Tên nhóm cửa hàng (phục vụ phân tầng và stress-test) \\
    \code{source}            & Nguồn dữ liệu (\code{mc\_ocr\_2021} hoặc \code{self\_collected}) \\
    \code{image\_path}       & Đường dẫn tương đối đến tệp ảnh \\
    \code{width}             & Chiều rộng ảnh gốc (pixel) \\
    \code{height}            & Chiều cao ảnh gốc (pixel) \\
    \code{annotation\_level} & Mức độ chú thích (\code{json\_and\_boxes} hoặc \code{json\_only}) \\
    \code{raw\_target}       & Nhãn thô ban đầu trước chuẩn hóa (JSON object) \\
    \code{target}            & Nhãn sạch đã qua chuẩn hóa (JSON object) \\
    \code{field\_boxes}      & Tọa độ bounding box $[x_1, y_1, x_2, y_2]$ của từng trường \\
    \code{field\_polygons}   & Tọa độ đa giác của từng trường thông tin \\
    \code{ocr\_cache\_path}  & Đường dẫn đến tệp lưu kết quả OCR cache \\
    \bottomrule
  \end{tabularx}
\end{table}

Ví dụ một dòng dữ liệu trong tệp JSONL thống nhất:
\begin{lstlisting}[style=json, caption={Ví dụ một dòng dữ liệu JSONL}, label={lst:jsonl-example}]
{
  "id": "mcocr_0001",
  "group_id": "g_0001",
  "store_group": "circlek",
  "source": "mcocr",
  "image_path": "images/mcocr_0001.jpg",
  "width": 1080, "height": 1920,
  "annotation_level": "json_and_boxes",
  "raw_target": {
    "store_name": "CIRCLE K",
    "address": "123 Nguyen Trai, HN",
    "date": "15/03/2021",
    "total": "115.000"
  },
  "target": {
    "store_name": "circle k",
    "address": "123 nguyen trai, hn",
    "date": "2021-03-15",
    "total": "115000"
  },
  "field_boxes": {"store_name": [100, 50, 400, 90]},
  "field_polygons": {},
  "ocr_cache_path": "ocr_cache/mcocr_0001.json"
}
\end{lstlisting}

% ============================================================================
\subsection{Chiến lược chia dữ liệu (Data Split) và ngăn ngừa Leakage}
\label{subsec:data-split}

\paragraph{Rủi ro rò rỉ dữ liệu (Data Leakage)}
Trong bài toán trích xuất thông tin biên lai, hiện tượng rò rỉ dữ liệu xảy ra khi các hình ảnh chụp cùng một biên lai vật lý (nhưng ở các góc máy khác nhau) bị phân chia nằm ở cả hai tập huấn luyện (Train) và tập kiểm thử (Test). Điều này làm cho mô hình dễ dàng "ghi nhớ" bố cục và nội dung cụ thể, dẫn đến kết quả đánh giá bị sai lệch và không phản ánh đúng khả năng tổng quát hóa thực tế của mô hình.

\paragraph{Biện pháp giải quyết trong Đồ án}
Để loại bỏ triệt để nguy cơ rò rỉ dữ liệu, nhóm nghiên cứu thiết lập chiến lược phân chia Main Split tỷ lệ \textbf{70/15/15} (Train/Val/Test) với hai ràng buộc kỹ thuật nghiêm ngặt:
1. \textbf{Group-based Split:} Việc phân chia được thực hiện ở cấp độ `group_id` (chứ không phải cấp tệp ảnh đơn lẻ). Toàn bộ ảnh có chung `group_id` bắt buộc phải nằm trong cùng một tập (hoặc Train, hoặc Val, hoặc Test).
2. \textbf{MD5 Hash Check:} Việc kiểm tra mã băm MD5 đảm bảo không có bất kỳ ảnh nào giống hệt nhau tồn tại song song giữa các tập split.
3. \textbf{Phân tầng theo nguồn (Stratified by source):} Đảm bảo tỷ lệ phân bố giữa dữ liệu MC-OCR và dữ liệu tự chụp là tương đương giữa các tập huấn luyện, phát triển và kiểm thử.

Thống kê số lượng mẫu sau phân chia được trình bày trong \cref{tab:split-stats}.

\begin{table}[htbp]
  \centering
  \caption{Phân bố số lượng mẫu sạch theo các tập dữ liệu.}
  \label{tab:split-stats}
  \begin{tabular}{lccc}
    \toprule
    \textbf{Tập dữ liệu} & \textbf{MC-OCR 2021} & \textbf{Tự thu thập} & \textbf{Tổng số mẫu} \\
    \midrule
    Train (Huấn luyện) & 766 & 325 & 1.091 \\
    Val (Phát triển)   & 164 & 70  & 234   \\
    Test (Kiểm thử)    & 164 & 70  & 234   \\
    \midrule
    \textbf{Tổng} & \textbf{1.094} & \textbf{465} & \textbf{1.559} \\
    \bottomrule
  \end{tabular}
\end{table}

\paragraph{Thống kê tỷ lệ thiếu nhãn trên tập kiểm thử (Test Set)}
Trong thực tế, một số biên lai bị thiếu thông tin hoặc thông tin bị mờ không thể đọc được. \cref{tab:missing-rate} thống kê số mẫu thiếu nhãn (missing rate) của từng trường cụ thể trên tập kiểm thử (Test set) gồm \textbf{234 mẫu}, phản ánh độ thử thách của tập kiểm định thực tế.

\begin{table}[htbp]
  \centering
  \caption{Tỷ lệ thiếu nhãn theo từng trường trên tập kiểm thử (234 mẫu).}
  \label{tab:missing-rate}
  \begin{tabular}{lcc}
    \toprule
    \textbf{Trường thông tin} & \textbf{Số mẫu thiếu nhãn} & \textbf{Tỷ lệ thiếu (\%)} \\
    \midrule
    \field{store\_name} & 10 & 4,27\% \\
    \field{date}        & 18 & 7,69\% \\
    \field{total}       & 17 & 7,26\% \\
    \field{address}     & 11 & 4,70\% \\
    \bottomrule
  \end{tabular}
\end{table}

% ============================================================================
\subsection{Kiểm tra tính hợp lệ dữ liệu (Validation)}
\label{subsec:data-validation}

Trước khi đưa vào huấn luyện, toàn bộ 1.559 mẫu dữ liệu sạch được chạy qua script tự động \code{validate\_dataset} để xác minh cú pháp và định dạng nhãn. Script thực hiện kiểm tra:
\begin{itemize}
  \item Định dạng ngày \field{date} phải tuân thủ chuẩn ISO \code{YYYY-MM-DD}.
  \item Giá trị trường \field{total} phải là chuỗi số nguyên thuần túy (không chứa ký tự chữ cái, đơn vị tiền tệ hoặc dấu phân cách hàng nghìn).
  \item Đường dẫn \code{image\_path} phải hợp lệ và tệp ảnh không bị lỗi đọc ghi (corrupted image).
\end{itemize}

Kết quả chạy thử nghiệm ghi nhận 15 mẫu có lỗi định dạng ngày tháng (chứa ký tự giờ phút giây ở nhãn chuẩn hóa) và 8 mẫu có trường \field{total} chứa dấu phân cách hàng nghìn. Nhóm nghiên cứu đã khắc phục triệt để các lỗi này bằng cách tinh chỉnh các bộ lọc chuẩn hóa văn bản tự động, đưa tỷ lệ lỗi hợp lệ về 0\% trước khi chính thức huấn luyện.
